\SourceLeaf{371}%
that the certainty of geometry has no basis in intuition and rests on definitions and on the principles of identity and contradiction. \textsc{Kant} presumably did not mean to deny that these logical aids are constantly used to give form to the truths of geometry and link them together. But surely no one familiar with the nature of geometry can fail to see that by themselves they achieve nothing and produce only barren flowers unless the fertilizing influence of a living intuition of the subject is present throughout. Mr.~\textsc{Schwab}'s objection, moreover, seems in part to rest on a misunderstanding. At least, judging from Section~16 of his work, which calls on the faculty of intuition from beginning to end and is meant to end by proving postulata \textsc{Euclidis} in generaliora resolvi posse, non sensu et \textit{intuitione} sed \textit{intellectu} fundata, it seems to us that Mr.~\textsc{Schwab} must have understood these names for different branches of the faculty of cognition differently from the Königsberg philosopher.

Although the author of the second work has treated his subject in an altogether different and genuinely mathematical way, we can judge the result no more favourably. We do not intend to set out the whole course of his attempted proof here. We will simply single out the point on which everything essentially depends. Imagine two straight lines intersecting at right angles at \(N\), and from a point \(S\) lying off those lines but in the same plane drop perpendiculars to them, \(ST\) and \(SM\). The point now is to prove that \(MST\) is a right angle. The author attempts a proof by contradiction. He first assumes that \(MST\) is acute, drops from \(T\) to \(MS\) the perpendicular \(Tp\), and proves that \(p\) must fall between \(S\) and \(M\). He then drops from \(p\) to \(NT\) the perpendicular \(pq\), so that \(q\) will fall between \(T\) and \(N\). He next drops from \(q\) to \(MS\) the perpendicular \(qp'\), so that \(p'\) will lie between \(p\) and \(M\). He then drops from \(p'\) to \(NT\) the perpendicular \(p'q'\), and so on. These operations can be continued without end, and so from the line \(MS\) the pieces \(Sp,pp'\) and so on are cut off one after another, each with an assignable magnitude and with no limit to their number. The author now takes this to be a contradiction, since \(MS\) would necessarily have to be exhausted in this way in the end.
